\documentclass[11pt]{article}

\usepackage[T1]{fontenc}
\usepackage[utf8]{inputenc}
\usepackage{lmodern}
\usepackage[margin=1.1in]{geometry}
\usepackage{amsmath,amssymb,amsthm,mathtools}
\usepackage{microtype}
\usepackage{xcolor}
\usepackage{array}
\usepackage{enumitem}
\usepackage{titlesec}
\usepackage{hyperref}

\titleformat{\section}
  {\normalfont\LARGE\bfseries}{\thesection}{1em}{}
\titleformat{\subsection}
  {\normalfont\Large\bfseries}{\thesubsection}{1em}{}

\hypersetup{
  pdftitle={Finite-Etale Span Projectors and Marked Galois-Sheet Blocks in Lean},
  pdfauthor={Chloe van der Vlugt},
  colorlinks=true,
  linkcolor=blue!55!black,
  citecolor=green!40!black,
  urlcolor=blue!60!black
}

\newtheorem{theorem}{Theorem}[section]
\newtheorem{proposition}[theorem]{Proposition}
\newtheorem{corollary}[theorem]{Corollary}
\theoremstyle{definition}
\newtheorem{definition}[theorem]{Definition}
\theoremstyle{remark}
\newtheorem*{remark}{Remark}

\DeclareMathOperator{\Corr}{Corr}
\DeclareMathOperator{\End}{End}
\DeclareMathOperator{\ev}{ev}
\DeclareMathOperator{\Gal}{Gal}
\DeclareMathOperator{\CH}{CH}
\DeclareMathOperator{\Def}{Def}

\newcommand{\Q}{\mathbf Q}
\newcommand{\one}{\mathbf 1}
\newcommand{\fe}{\mathrm{f\acute et}}
\newcommand{\Geta}{\Gamma_{\eta}}
\newcommand{\Gbar}{\overline{\Gamma}}
\newcommand{\leanpath}[1]{\texttt{\nolinkurl{#1}}}
\newcommand{\corrsrc}[1]{%
  \leanpath{HodgeConjecture/Correspondence/#1}}
\newcommand{\sheetsrc}[1]{%
  \leanpath{HodgeConjecture/Boundary/Sheets/#1}}
\newcommand{\normsrc}[1]{%
  \leanpath{HodgeConjecture/Boundary/Normalization/#1}}
\newcommand{\boundarysrc}[1]{%
  \leanpath{HodgeConjecture/Boundary/#1}}
\newcommand{\chowsrc}[1]{%
  \leanpath{HodgeConjecture/Realization/Chow/#1}}

\newenvironment{leanfiles}
  {\par\medskip\begingroup\footnotesize\raggedright\noindent
   \textit{Lean correspondence.}\ }
  {\par\endgroup\medskip}

\title{Finite-\'Etale Span Projectors and\\
  Marked Galois-Sheet Blocks in Lean}
\author{Chloe van der Vlugt}
\date{Draft of September 11, 2026}

\begin{document}

\maketitle

\begin{abstract}
We formalize in Lean~4 a convolution algebra built from isomorphism classes
of scheme endocorrespondences whose legs are finite \'etale.  An action by
deck transformations induces a rational algebra homomorphism from the group
algebra to this free correspondence algebra; consequently, idempotents and
finite ordered orthogonal decompositions transport to formal correspondence
projectors and diagonal decompositions.  For finite groups we obtain
normalized averaging and complementary projectors, and for finite subgroups
we construct the normalized Hecke corner and its double-coset operators.  For
a supplied comparison square between generic and compactified graph
realizations, every group-algebra idempotent has a compactified idempotent
image with the prescribed generic restriction.  For an arbitrary raw closure
we prove only that its convolution-idempotence defect restricts to zero, and
isolate zero defect as the additional condition recovering an idempotent
extension.  For
a marked normal closure, we also construct the induced deck action on the
affine integral-closure normalization and identify normalization in a fixed
field with the corresponding fixed subring.  Canonical conjugate
normalization legs then define fixed--moving boundary ideals; under explicit
center, ramification, and boundary-separation hypotheses, they rule out
codimension-one ramification.

For a marked quotient \(G/H\), we prove that multiplication by the base-sheet
selector in \(\Q^{G/H}\) is \(\operatorname{diag}(1,0)\) in selector
coordinates and the graph projector of \(a_H=q_H-1\) in constant-unit coordinates.
The two forms are related by an explicit shear, and the comparison is natural
under rational maps of the residual factor.  We finally isolate a conditional
realization criterion: a supplied extension package and chosen compactified
idempotent, residual realization, action-intertwining square, and lower-left
detection map place the resulting class in the range of a supplied
cycle-class model.
We also record a forward conditional boundary corollary: different avoidance
produces a prescribed compactified-algebra idempotent only after cycle
extraction, compactified idempotence, one-way purity, inertia, restriction,
and prescribed-action inputs are separately supplied.

The free span algebra constructed here is not a Chow correspondence algebra.
No concrete Chow realization, boundary extension, or new case of the Hodge
conjecture is proved.
\end{abstract}

\begin{center}
  \small Manuscript draft for author review.  No license assertion is made
  in this draft.  A literature and novelty review is still pending.
\end{center}

\begingroup
\small
\tableofcontents
\endgroup

\section{Introduction}
\label{sec:introduction}

Let \(p_X:X\to S\) be a morphism of schemes.  A correspondence from \(X\)
to itself over \(S\) is a span

\[
  X \xleftarrow{\ell} Z \xrightarrow{r} X
\]

whose two maps to \(S\) agree.  Two such spans compose by taking the fiber
product over the middle copy of \(X\).  The diagonal span is the identity.
Raw fiber-product composition is associative and unital only up to canonical
isomorphism; passage to carrier-isomorphism classes makes these laws strict.

The present note studies the free rational algebra on the classes whose two
legs are finite \'etale.  This setting is large enough to contain graphs of
automorphisms, is closed under composition and transpose, and is small enough
to avoid imposing additivity on arbitrary infinite coproducts.  Its product
is convolution of correspondences.

Three algebras must be kept distinct throughout:

\begin{enumerate}[label=\textup{(\roman*)},leftmargin=2.6em]
  \item the pointwise function algebra on a finite set of generic sheets;
  \item the group algebra, and its Hecke corners, with convolution product;
  \item a prospective Chow correspondence algebra on a compactification.
\end{enumerate}

A characteristic function in the first algebra is pointwise idempotent.  That
fact alone does not say that a sum of graph correspondences is idempotent in
the second or third algebra.  For example, the graph of a nonidentity
automorphism \(g\) satisfies

\[
  [\Gamma_g]^2=[\Gamma_{g^2}],
\]

not \([\Gamma_g]^2=[\Gamma_g]\) in general.  The convolution projectors in
this note instead come from actual idempotents of a group algebra or a Hecke
corner.  The marked function-algebra selector enters later through a separate
operator comparison.

The two principal formal results are as follows.

\begin{itemize}[leftmargin=2.2em]
  \item A deck action induces an algebra map from the rational group algebra
        to the finite-\'etale free correspondence algebra.  It transports
        group-algebra idempotents and orthogonal decompositions to
        convolution projectors and diagonal decompositions.
  \item A supplied generic--compactified graph square transports those
        idempotents to a supplied compactified algebra with zero convolution
        defect.  For arbitrary closure classes, restriction kills the defect
        but does not force it to vanish.
  \item Multiplication by a marked sheet selector is diagonal in its
        image--kernel coordinates and graph-shaped in the splitting marked
        by the constant unit.  The transition between these descriptions is
        an explicit lower shear.
\end{itemize}

Section~\ref{sec:conditional-realization} gives a conditional interface
connecting the second result to a supplied cycle-class model.  The interface
exposes rather than proves the missing geometry: compactification, a residual
realization, an intertwining square, and lower-left block detection.
The same section records the honest current form of the boundary slogan as a
forward conditional corollary with every remaining geometric input visible.

\paragraph{Literature status.}
The underlying group-algebra and graph calculations are classical.  This
draft is positioned only as a formalization and infrastructure report.  A
systematic comparison with prior formalizations and the mathematical
literature has not yet been completed, so no novelty or priority claim is
made here.

\subsection*{Semantic notation}

We reserve one symbol for each mathematical role.  In particular, symbols
are not shared merely because two Lean structures have similarly named
fields.  The base, marked extension, and normal closure are

\[
  B\subset K=\operatorname{Frac}(B),\qquad
  L/K,\qquad N/K,\qquad
  G=\Gal(N/K),\qquad H=\Gal(N/\iota(L)).
\]

Their normalization rings and schemes are

\[
  \widetilde B_L=\operatorname{IntCl}_B(L),\quad
  \widetilde B_N=\operatorname{IntCl}_B(N),\qquad
  Y_L=\operatorname{Spec}(\widetilde B_L),\quad
  Y_N=\operatorname{Spec}(\widetilde B_N).
\]

The conjugate embedding and normalization leg are \(\iota_g\) and
\(c_g:Y_N\to Y_L\).  Their pulled-back boundary ideal is \(J_g\); for
\(I\leq G\), the fixed, moving, and combined ideals are respectively
\(J_{\mathrm{fix}}(I)\), \(J_{\mathrm{mov}}(I)\), and
\(J_{\mathrm{fm}}(I)\).  We write
\(\operatorname{Ram}^{(1)}(Y_N/B)\) for the ramified height-one locus and
\(\operatorname{Sep}_{\partial}(\mathcal C)\) for conjugate-boundary
separation.

For a marked sheet set \(\Sigma_H=G/H\), the pointwise algebra is
\(F_k(H)=k^{G/H}\), a selected-set characteristic function is \(q_A\), and
the base-sheet selector is \(q_H=\one_{\{H\}}\).  By contrast,

\[
  e_H=|H|^{-1}\sum_{h\in H}[h]\in\Q[G]
\]

always denotes the group-algebra convolution idempotent, while
\(\pi_H^\eta=\gamma_\eta(e_H)\) denotes its generic graph-correspondence
image.  We use \(\Geta\) and \(\Gbar\)
for the generic and compactified correspondence algebras and
\(r:\Gbar\to\Geta\) for restriction.  A proposed compactified projector is
\(\Pi\), and its unique idempotence-defect notation is

\[
  \Def(\Pi)=\Pi^2-\Pi.
\]

Finally, \(E_H=\ker(\ev_H)\) is the marked residual sheet module,
\(a_H=q_H-1\), \(P_\alpha(t,x)=(t,t\alpha)\) is the graph projector, and
\(U_\alpha(t,x)=(t,x+t\alpha)\) is its shear.  The compactified action is
\(\rho:\Gbar\to\End_{\Q}(\Q\oplus V)\), while \(\mathsf{Ch}\) and
\(\operatorname{cl}:\mathsf{Ch}\to V\) denote the independently supplied
cycle module and cycle-class map.  Thus \(q_A\), \(e_H\), and \(\Pi\) live
in three different algebras and are never identified without an explicit
comparison map.

\section{Scheme correspondences and the free convolution algebra}
\label{sec:correspondences}

Fix morphisms \(p_X:X\to S\) and \(p_Y:Y\to S\).  A relative scheme
correspondence \(P:X\dashrightarrow Y\) consists of a carrier \(P_0\), two
maps

\[
  X\xleftarrow{\,P_\ell\,}P_0
   \xrightarrow{\,P_r\,}Y,
  \qquad
  P_\ell p_X=P_r p_Y,
\]

where compositions are written in diagrammatic order in this display.  The
compatibility supplies a canonical support map

\[
  P_0\longrightarrow X\times_S Y.
\]

If \(P:X\dashrightarrow Y\) and \(Q:Y\dashrightarrow T\), their composite
\(Q\circ P\) has carrier \(P_0\times_YQ_0\), left map induced by \(P_\ell\),
and right map induced by \(Q_r\).  Canonical pullback isomorphisms give

\[
  1_Y\circ P\simeq P,
  \qquad
  P\circ1_X\simeq P,
  \qquad
  R\circ(Q\circ P)\simeq(R\circ Q)\circ P.
\]

An isomorphism of correspondences is an isomorphism of carriers commuting
with both legs.  Write \(\langle P\rangle\) for the corresponding isomorphism
class.

\begin{proposition}[Strict convolution after quotienting]
\label{prop:class-monoid}
For fixed \(p_X:X\to S\), isomorphism classes of endocorrespondences form a
monoid under

\[
  \langle Q\rangle\,\langle P\rangle
  =\langle Q\circ P\rangle.
\]

Its unit is the class of the diagonal correspondence.
\end{proposition}

\begin{proof}
Composition respects carrier isomorphisms.  The canonical unit and
associativity isomorphisms above become equalities after quotienting, which
are precisely the monoid laws.
\end{proof}

\begin{leanfiles}
The raw span, composition, carrier isomorphism, unit isomorphisms, and
associativity isomorphism are in \corrsrc{Scheme.lean}.  The quotient and
monoid instance are in \corrsrc{Classes.lean}, under
\leanpath{SchemeCorrespondence.Class} and
\leanpath{SchemeCorrespondence.classMonoid}.
\end{leanfiles}

Call \(P\) finite \'etale when both \(P_\ell\) and \(P_r\) are finite
\'etale.  This condition is invariant under correspondence isomorphism and
is preserved by the diagonal, transpose, and fiber-product composition.  It
therefore cuts out a submonoid

\[
  \Corr_{\fe}(X/S)
  \subseteq \Corr(X/S).
\]

We define the finite-\'etale free correspondence algebra by

\begin{equation}
\label{eq:free-algebra}
  \mathcal C_{\fe}^{\mathrm{free}}(X/S)
  :=\Q[\Corr_{\fe}(X/S)].
\end{equation}

Thus a finite-\'etale span class is a basis vector, multiplication is induced
by fiber-product composition, and the diagonal class is the unit.

\begin{theorem}[Finite-\'etale correspondence algebra]
\label{thm:finite-etale-algebra}
Finite-\'etale correspondence classes are closed under convolution and form
a monoid.  The algebra \eqref{eq:free-algebra} is consequently a unital,
generally noncommutative rational algebra.  Transpose extends linearly to an
involution satisfying

\[
  (ab)^{\mathsf t}=b^{\mathsf t}a^{\mathsf t},
  \qquad
  (a^{\mathsf t})^{\mathsf t}=a.
\]
\end{theorem}

\begin{proof}
Finiteness and \'etaleness are stable under isomorphism, composition, and
base change.  For a composite, each new leg is a base-changed old leg
followed by the other old leg.  This proves closure.  Pullback symmetry
identifies the transpose of \(Q\circ P\) with
\(P^{\mathsf t}\circ Q^{\mathsf t}\); quotienting and extending linearly give
the displayed anti-multiplicative involution.
\end{proof}

\begin{leanfiles}
Closure of the admissible class is in \corrsrc{Etale.lean} and
\corrsrc{EtaleClasses.lean}.  The free algebra is
\leanpath{FiniteEtaleFreeAlgebra}.  Transpose and its ring involution are in
\corrsrc{TransposeAlgebra.lean}; the principal multiplication theorem is
\leanpath{finiteEtaleFreeTranspose_mul}.
\end{leanfiles}

\begin{remark}[Not a Chow algebra]
The construction is intentionally free.  It does not impose rational
equivalence, pure-dimension conditions, or refined-intersection relations.
It also does not identify a disjoint union with a sum:

\[
  \langle P\amalg Q\rangle
  =\langle P\rangle+\langle Q\rangle
\]

is not a relation in
\(\mathcal C_{\fe}^{\mathrm{free}}(X/S)\).  Such an equality belongs to a
separately supplied additive realization.
\end{remark}

\section{Deck actions and rational projectors}
\label{sec:deck-projectors}

An automorphism of \(X\) over \(S\) has graph correspondence

\[
  \Gamma_\sigma=
  \bigl(X\xleftarrow{1_X}X\xrightarrow{\sigma}X\bigr).
\]

Both legs are isomorphisms and hence finite \'etale.  Under the convention of
Proposition~\ref{prop:class-monoid}, graph composition agrees with
automorphism multiplication:

\begin{equation}
\label{eq:graph-product}
  \langle\Gamma_\sigma\rangle
  \langle\Gamma_\tau\rangle
  =\langle\Gamma_{\sigma\tau}\rangle.
\end{equation}

Consequently an action

\[
  a:G\longrightarrow\operatorname{Aut}_S(X)
\]

induces an algebra homomorphism

\begin{equation}
\label{eq:deck-algebra-map}
  \gamma_a:\Q[G]\longrightarrow
  \mathcal C_{\fe}^{\mathrm{free}}(X/S),
  \qquad
  [g]\longmapsto\langle\Gamma_{a(g)}\rangle.
\end{equation}

The Lean name \leanpath{FiniteDeckAction} denotes an action by a group of
deck transformations.  Despite the name, the definition itself neither
requires \(G\) to be finite nor asserts that \(p_X\) is a finite-\'etale
cover.

\begin{theorem}[Deck-action projector transport]
\label{thm:deck-transport}
Let \(p_X:X\to S\) be a scheme morphism and let \(G\) act on \(X\) by
\(S\)-automorphisms.  The graph assignment induces the rational algebra map
\eqref{eq:deck-algebra-map}.  It has the following consequences.

\begin{enumerate}[label=\textup{(\roman*)},leftmargin=2.6em]
  \item If \(e\in\Q[G]\) satisfies \(e^2=e\), then
        \(\gamma_a(e)^2=\gamma_a(e)\) for correspondence convolution.
  \item Let \(I\) be finite and let \((e_i)_{i\in I}\) be idempotents in
        \(\Q[G]\) satisfying
        \[
          e_ie_j=0\quad(i\ne j),
          \qquad
          \sum_{i\in I}e_i=1.
        \]
        Then \((\gamma_a(e_i))_{i\in I}\) is an ordered orthogonal family of
        convolution idempotents whose sum is the diagonal unit.
  \item If \(G\) is finite, the normalized average
        \[
          e_G=\frac{1}{|G|}\sum_{g\in G}[g]
        \]
        and \(1-e_G\) map to complementary convolution idempotents.
\end{enumerate}
\end{theorem}

\begin{proof}
Equation~\eqref{eq:graph-product} and the identity-graph comparison produce
the algebra homomorphism.  Algebra homomorphisms preserve multiplication,
zero, one, addition, and finite sums, so they preserve the equations in
parts~(i) and~(ii).  For part~(iii), multiplication by any basis element
permutes the finite group sum.  Hence

\[
  \left(\sum_{g\in G}[g]\right)^2
  =|G|\sum_{g\in G}[g],
\]

and \(e_G^2=e_G\).  The complement of any idempotent is idempotent and is
orthogonal to it.
\end{proof}

\begin{leanfiles}
Graph composition and the graph algebra map are in \corrsrc{Graph.lean}.
The action map \leanpath{FiniteDeckAction.toFiniteEtaleAlgebra}, the average,
and its complement are in \corrsrc{FiniteGroupProjector.lean}.  Transport of
arbitrary idempotents and finite ordered orthogonal decompositions is in
\corrsrc{GroupAlgebraProjector.lean}, principally
\leanpath{groupAlgebraCorrespondence_isIdempotent} and
\leanpath{GroupAlgebraDiagonalDecomposition.correspondence_sum_eq_diagonal}.
\end{leanfiles}

\begin{remark}[Centrality]
Centrality of \(e\) is not required for Theorem~\ref{thm:deck-transport}.
Central idempotents are useful when one wants canonical isotypic summands,
but idempotence transport uses only \(e^2=e\).
\end{remark}

\section{Compactified graph comparison and closure defects}
\label{sec:compactified-graph}

The group-algebra calculation also separates a useful sufficient condition
for compactified idempotence.  Let
\(p_\eta:X_\eta\to S_\eta\) and
\(\bar p:\bar X\to\bar S\) carry actions \(a_\eta\) and \(\bar a\) of the
same group \(G\).  Supply rational algebras \(\Geta\) and \(\Gbar\),
algebra maps

\[
  c_\eta:\mathcal C_{\fe}^{\mathrm{free}}(X_\eta/S_\eta)
    \longrightarrow \Geta,
  \qquad
  \bar c:\mathcal C_{\fe}^{\mathrm{free}}(\bar X/\bar S)
    \longrightarrow\Gbar,
\]

and a restriction map \(r:\Gbar\to\Geta\).  The comparison datum is the
commutative square on the rational group algebra

\begin{equation}
\label{eq:compactified-graph-square}
  r\,\bar c\,\gamma_{\bar a}
  =c_\eta\,\gamma_{a_\eta}:\Q[G]\longrightarrow\Geta.
\end{equation}

The word ``compactified'' names the intended application.  The formal datum
does not itself relate the two schemes, construct either cycle map, or
identify \(\Gbar\) with a Chow ring.

\begin{proposition}[Projectors through a supplied compactified graph square]
\label{prop:compactified-graph-projector}
If \(e\in\Q[G]\) satisfies \(e^2=e\), put

\[
  \Pi_e=\bar c\bigl(\gamma_{\bar a}(e)\bigr)\in\Gbar.
\]

Then

\[
  \Pi_e^2=\Pi_e,
  \qquad
  r(\Pi_e)=c_\eta\bigl(\gamma_{a_\eta}(e)\bigr),
  \qquad
  \Pi_e^2-\Pi_e=0.
\]

The same construction for the normalized average of a finite subgroup
produces a compactified projector-extension package with a singleton witness
type.
\end{proposition}

\begin{proof}
Both graph maps and both supplied cycle maps are algebra homomorphisms, so
they preserve the equation \(e^2=e\).  The restriction identity is exactly
Equation~\eqref{eq:compactified-graph-square} evaluated at \(e\).  The defect
identity is the idempotence equation rewritten by subtraction.
\end{proof}

There is a weaker construction for a closure not known to arise from a
global graph map.  Let \(e_\eta\in\Geta\) be idempotent and let a raw closure
class \(\Pi\in\Gbar\) satisfy \(r(\Pi)=e_\eta\).  Define its convolution
defect by

\[
  \Def(\Pi)=\Pi^2-\Pi.
\]

Then algebra-map compatibility gives only

\begin{equation}
\label{eq:defect-restricts-zero}
  r\bigl(\Def(\Pi)\bigr)
  =e_\eta^2-e_\eta=0.
\end{equation}

Thus the defect belongs to the kernel of generic restriction.  The subtype
of closure witnesses satisfying \(\Def(\Pi)=0\) canonically forms an
\leanpath{AlgebraProjectorExtension}.

\begin{leanfiles}
The comparison square, compactified projector, restriction and idempotence
theorems, zero-defect theorem, and subgroup-average specialization are in
\boundarysrc{CompactifiedGraphProjector.lean}, under
\leanpath{CompactifiedDeckGraphComparison}.  Raw closure classes and
Equation~\eqref{eq:defect-restricts-zero} are in
\boundarysrc{ProjectorClosureDefect.lean}, under
\leanpath{AlgebraProjectorClosure}.
\end{leanfiles}

\begin{remark}[Kernel membership is not boundary vanishing]
Equation~\eqref{eq:defect-restricts-zero} is the precise formal precursor to
a boundary-support statement.  The present model neither identifies the
restriction kernel with boundary-supported Chow classes nor proves that a
class in that kernel vanishes.  Proposition
~\ref{prop:compactified-graph-projector} avoids this issue only because the
supplied global graph algebra map makes idempotence formal from the outset.
\end{remark}

\section{Finite-subgroup Hecke corners}
\label{sec:hecke}

Let \(H\leq G\) be finite, with \(G\) not necessarily finite, and put

\[
  e_H=\frac{1}{|H|}\sum_{h\in H}[h]\in\Q[G].
\]

The element \(e_H\) is idempotent.  Its corner

\[
  \mathcal H(G,H)=e_H\Q[G]e_H
\]

is a unital ring whose unit is \(e_H\), not the ambient unit of
\(\Q[G]\).  For \(g\in G\), define the normalized sandwiched operator

\[
  T_g=e_H[g]e_H\in\mathcal H(G,H).
\]

\begin{theorem}[Hecke-corner transport]
\label{thm:hecke}
The element \(T_g\) depends only on the double coset \(HgH\).  The deck-action
map \(\gamma_a\) restricts to a unital ring homomorphism

\[
  e_H\Q[G]e_H
  \longrightarrow
  \gamma_a(e_H)
    \mathcal C_{\fe}^{\mathrm{free}}(X/S)
    \gamma_a(e_H).
\]

Every idempotent in the source corner therefore determines a convolution
idempotent in the finite-\'etale correspondence algebra.
\end{theorem}

\begin{proof}
For \(h,k\in H\), left and right invariance of the subgroup average give

\[
  e_H[hgk]e_H=e_H[g]e_H.
\]

This proves double-coset invariance.  A unital ring homomorphism sends an
idempotent corner to the corner of its image; applying this construction to
\(\gamma_a\) gives the displayed homomorphism and preserves idempotence.
\end{proof}

\begin{leanfiles}
The general corner map, normalized subgroup average, double-coset
factorization, and correspondence-corner map are in
\corrsrc{HeckeProjector.lean}.  The key declarations are
\leanpath{SubgroupHeckeAlgebra}, \leanpath{heckeElement_eq_of_rel},
\leanpath{heckeToCorrespondenceCorner}, and
\leanpath{heckeCorrespondence_isIdempotent}.
\end{leanfiles}

\begin{remark}[Normalization and scope]
The element \(T_g=e_H[g]e_H\) is not the raw characteristic sum of the
finite set \(HgH\); the two differ by a stabilizer-dependent rational factor.
The construction neither proves that every \(T_g\) is idempotent nor asserts
commutativity, a double-coset basis, or structure constants for the corner.
\end{remark}

\section{The conditional Galois collision formula}
\label{sec:galois-collision}

Let \(f:X\to Y\) be finite \'etale and let \(G\) be finite.  The formalized
Galois collision statement does not infer a torsor decomposition merely from
these hypotheses.  Instead, it accepts explicit data consisting of

\begin{enumerate}[label=\textup{(\roman*)},leftmargin=2.6em]
  \item an action of \(G\) on \(X\) over \(Y\);
  \item an isomorphism
        \[
          \coprod_{g\in G}X\xrightarrow{\sim}X\times_YX;
        \]
  \item compatibility of the two collision projections with the identity
        and \(g\) on the \(g\)-summand.
\end{enumerate}

These fields comprise \leanpath{GaloisCollisionDecomposition}.  To express
finite additivity, let \(Y\to S\) be a further base map and let

\[
  \chi:\mathcal C_{\fe}^{\mathrm{free}}(X/S)\longrightarrow\Gamma
\]

be a rational algebra homomorphism which, as separately supplied data, maps
finite admissible coproducts to sums.  This is the interface
\leanpath{FiniteEtaleAlgebraModel}.

\begin{theorem}[Galois collision comparison, given the torsor data]
\label{thm:galois-collision}
Under the data above,

\begin{equation}
\label{eq:galois-graph-sum}
  \chi\bigl([X\times_YX]\bigr)
  =\sum_{g\in G}\chi([\Gamma_g]).
\end{equation}

Consequently the normalized collision class

\[
  c_G=\frac{1}{|G|}\chi\bigl([X\times_YX]\bigr)
\]

is an idempotent in \(\Gamma\).
\end{theorem}

\begin{proof}
The supplied torsor isomorphism identifies the collision correspondence with
the indexed coproduct of the graph correspondences.  Finite additivity of
\(\chi\) gives \eqref{eq:galois-graph-sum}.  The right side, after division
by \(|G|\), is the image under \(\chi\) of the graph average from
Theorem~\ref{thm:deck-transport}; it is therefore idempotent.
\end{proof}

\begin{leanfiles}
The decomposition record, graph-sum theorem, comparison with the group
average, and idempotence theorem are in \corrsrc{GaloisCollision.lean}:
\leanpath{GaloisCollisionDecomposition},
\leanpath{classOf_collision_eq_sum_graph},
\leanpath{averagingProjector_realizes_normalizedCollisionClass}, and
\leanpath{normalizedCollisionClass_isIdempotent}.  The finite-additivity
interface is in \corrsrc{FiniteEtaleModel.lean}.
\end{leanfiles}

\begin{remark}[The off-diagonal class]
Under the same supplied Galois collision decomposition and model \(\chi\),
suppose one separately supplies a
\leanpath{CollisionSchemeDecomposition} \(E_{\mathrm{cl}}\) of \(f\).  Write
\(E_{\mathrm{off}}\) for its off-diagonal self-correspondence.  Additivity
then gives the realized identity

\[
  \chi([E_{\mathrm{off}}])
  =\sum_{g\in G}\chi([\Gamma_g])-1_{\Gamma}.
\]

This is an additive identity.  The off-diagonal class is not asserted to be
convolution-idempotent, normalized or otherwise.
\end{remark}

\section{Marked normal closures and generic sheets}
\label{sec:generic-sheets}

We next record how a marked field extension produces the sheet-function
algebra used in the block calculation.  Let \(L/K\) be finite and let \(N\)
be a supplied finite normal closure together with a distinguished embedding
\(\iota:L\hookrightarrow N\).  Put

\[
  G=\Gal(N/K),
  \qquad
  H=\{g\in G:g|_{\iota(L)}=1\}.
\]

Before constructing the generic evaluation, we record the compatible affine
normalization statement.  Let \(B\) be a commutative ring equipped with
compatible algebra maps to \(K\), \(L\), and \(N\).  Write
\(\operatorname{IntCl}_B(M)\) for the integral closure of \(B\) in a field
\(M\).

\begin{proposition}[Normal-closure deck action and fixed normalization]
\label{prop:fixed-normalization}
Every \(\sigma\in\Gal(N/K)\) restricts to a \(B\)-algebra automorphism of
\(\operatorname{IntCl}_B(N)\).  By contravariance of affine spectrum, these
automorphisms define a canonical action over \(\operatorname{Spec}B\) on
\[
  \operatorname{Spec}\operatorname{IntCl}_B(N).
\]
For every subgroup \(J\leq\Gal(N/K)\), there is a canonical \(B\)-algebra
equivalence
\[
  \operatorname{IntCl}_B(N^J)
  \simeq_B
  \bigl(\operatorname{IntCl}_B(N)\bigr)^J.
\]
If \(K\) is perfect and \(J=H\) is the subgroup fixing the marked copy of
\(L\), this specializes to
\[
  \operatorname{IntCl}_B(L)
  \simeq_B
  \bigl(\operatorname{IntCl}_B(N)\bigr)^H.
\]
\end{proposition}

\begin{proof}
A \(K\)-automorphism of \(N\) fixes the image of \(B\) because the maps
factor through \(K\).  Integrality over \(B\) is preserved by algebra
automorphisms, so restriction gives the first action.  Applying
\(\operatorname{Spec}\) to the inverse ring automorphism makes the scheme
construction multiplicative and preserves the structure map to
\(\operatorname{Spec}B\).

An element of \(\operatorname{IntCl}_B(N)\) is fixed by \(J\) exactly when
its image in \(N\) lies in \(N^J\).  Integrality over \(B\) is the same on
both sides, which gives the displayed fixed-ring equivalence.  For the marked
subgroup \(H\), finite Galois correspondence over a perfect base identifies
\(N^H\) with the distinguished image of \(L\), giving the specialization.
\end{proof}

\begin{leanfiles}
Restriction of automorphisms to the integral closure and the affine scheme
action are in \normsrc{DeckAction.lean}, under
\leanpath{normalClosureNormalizationRingEquiv},
\leanpath{normalClosureDeckTransformation}, and
\leanpath{normalClosureFiniteDeckAction}.  The fixed-ring equivalence and its
marked perfect-field specialization are
\leanpath{fixedFieldNormalizationEquivFixedRing} and
\leanpath{intermediateNormalizationEquivFixedRing} in
\normsrc{InvariantIntermediate.lean}.
\end{leanfiles}

\begin{remark}[No \'{e}taleness or quotient assertion]
Proposition~\ref{prop:fixed-normalization} constructs an action on the affine
normalization.  It does not assert that the normalization map is finite
\'etale, extend the action to an independently chosen compactification, or
prove a scheme-quotient universal property.
\end{remark}

Every \(g\in G\) also gives a conjugate \(B\)-algebra embedding
\(L\to N\), hence a canonical normalization leg

\[
  c_g:\operatorname{Spec}\operatorname{IntCl}_B(N)
    \longrightarrow
    \operatorname{Spec}\operatorname{IntCl}_B(L).
\]

It depends only on the right coset \(gH\) and lies over
\(\operatorname{Spec}B\).  If \(j:U\hookrightarrow
\operatorname{Spec}\operatorname{IntCl}_B(L)\) is a marked visible open, let
\(J_g\) be the vanishing ideal of the complement of \(c_g^{-1}(U)\).
For a subgroup \(I\leq G\), let \(I_{\mathrm{fix}}\) be generated by all
differences \(\sigma(t)-t\), with \(\sigma\in I\), and put

\[
  J_{\mathrm{fm}}(I)
  =I_{\mathrm{fix}}+
    \sum_{g:\,I\not\leq gHg^{-1}}J_g.
\]

The formal sum is implemented as a supremum of ideals, so no finiteness of
the indexing type is needed for its definition.

\begin{proposition}[Conjugate-boundary separation endpoint]
\label{prop:conjugate-boundary-separation}
Assume \(B\) is Noetherian and supply a normalization--inertia diagram whose
double-coset centers agree with the canonical legs \(c_g\), whose relative
inertia indices agree with geometric ramification indices.  Assume also that
the marked visible morphism to \(\operatorname{Spec}B\) is formally
unramified.  If

\[
  I\neq 1,\quad \operatorname{ht}(\mathfrak p)=1
  \quad\Longrightarrow\quad
  J_{\mathrm{fm}}(I)\not\subseteq\mathfrak p,
\]

for every subgroup \(I\leq G\) and every prime \(\mathfrak p\) of the
normal-closure normalization, then that normalization is unramified over
\(B\) at every height-one prime.
\end{proposition}

\begin{proof}
If a height-one prime \(\mathfrak p\) were ramified, its valuation inertia
group \(I_{\mathfrak p}\) would be nontrivial.  Residue-field triviality puts
every generator \(\sigma(t)-t\) in \(\mathfrak p\), so
\(I_{\mathrm{fix}}\subseteq\mathfrak p\).  Whenever
\(I_{\mathfrak p}\not\leq gHg^{-1}\), the corresponding relative inertia
index is nontrivial.  The supplied ramification comparison and formal
unramifiedness force that conjugate center into the deleted visible boundary;
the supplied center comparison then gives \(J_g\subseteq\mathfrak p\).
Consequently \(J_{\mathrm{fm}}(I_{\mathfrak p})\subseteq\mathfrak p\),
contradicting separation.
\end{proof}

\begin{corollary}[Conditional projector extension from separation]
\label{cor:separation-projector-extension}
Under the hypotheses of Proposition
~\ref{prop:conjugate-boundary-separation}, every selected set of generic
sheets satisfies all packaged height-one inertia-saturation conditions.
Consequently, a supplied selector-purity/gluing equivalence produces a
global selector extension.  A supplied subgroup-projector extension package
likewise produces a compactified correspondence projector.
\end{corollary}

\begin{proof}
The proposition makes the type of ramified height-one points empty, so every
local saturation condition follows by elimination.  The two conclusions
are then the forward directions of the supplied extension interfaces.
\end{proof}

\begin{leanfiles}
The conjugate embeddings, normalization legs, right-coset invariance, and
base compatibility are in \normsrc{ConjugateNormalization.lean}.  The
pulled-back boundary ideals, fixed--moving ideal, separation predicate, and
Proposition~\ref{prop:conjugate-boundary-separation} are in
\normsrc{ConjugateBoundary.lean}, under
\leanpath{ConjugateBoundarySeparation} and
\leanpath{noCodimensionOneRamification_of_conjugateBoundarySeparation}.
The selector and subgroup-projector wrappers in Corollary
~\ref{cor:separation-projector-extension} are in
\boundarysrc{SheetProjectorExtension.lean} and
\boundarysrc{GaloisProjectorExtension.lean}.
All four modules are native to this repository and import no sibling project.
\end{leanfiles}

\begin{remark}[The remaining purity arrow]
The proposition proves only the codimension-one conclusion.  It neither
proves the displayed separation hypothesis nor promotes that conclusion to
global \'etaleness.  The active formal development therefore stops at the
codimension-one conclusion and keeps the actual cover-extension and
correspondence-gluing statements as separately supplied inputs.
\end{remark}

The cosets \(G/H\) index the \(K\)-embeddings of \(L\) into \(N\).  Chosen
quotient representatives give an evaluation map

\begin{equation}
\label{eq:collision-sheet-evaluation}
  \epsilon:L\otimes_KL\longrightarrow (G/H\to N),
  \qquad
  \epsilon(x\otimes y)(gH)=\iota(x)\,g(\iota(y)).
\end{equation}

The formula is independent of the representative because \(H\) fixes the
marked copy of \(L\).

Assume \(L/K\) is formally unramified and essentially of finite type.  Its
chosen diagonal element

\[
  \Delta_{L/K}\in L\otimes_KL
\]

has multiplication image \(1\) and is killed by the usual differences
\(1\otimes x-x\otimes1\).  These properties identify its sheet evaluations.

\begin{proposition}[The diagonal element selects the marked sheet]
\label{prop:diagonal-selector}
Under the hypotheses above,

\[
  \epsilon(\Delta_{L/K})=q_H,
\]

where \(q_H\in (G/H\to N)\) is the characteristic function of the base coset
\(H\).  In particular, \(\Delta_{L/K}\) is idempotent.  If \(K\) is perfect,
normal-closure base change gives an \(N\)-algebra equivalence

\[
  N\otimes_KL\simeq (G/H\to N).
\]
\end{proposition}

\begin{proof}
At the base coset, evaluation is multiplication followed by the distinguished
embedding, so the value is \(1\).  At a nonbase coset, the associated
embedding differs from \(\iota\) on some \(x\in L\).  Applying the relation
\((1\otimes x-x\otimes1)\Delta_{L/K}=0\) and using that \(N\) is a field
forces the value to vanish.  Thus the image is \(q_H\).  Idempotence follows
from the same defining relations for the diagonal element.  Over a perfect
base, separability and the finite-dimensional rank comparison upgrade the
normal-closure evaluation map to the stated equivalence.
\end{proof}

\begin{leanfiles}
The marked normal-closure record and fixing subgroup are in
\normsrc{NormalClosure.lean}.  The evaluation maps and selector comparison
are in \normsrc{SelectorComparison.lean}, especially
\leanpath{NormalClosureData.collisionToGenericSheets_elem_eq_baseSheetSelector},
\leanpath{formallyUnramifiedElem_isIdempotent}, and
\leanpath{NormalClosureData.normalClosureTensorGenericSheetsEquiv}.
The intermediate integral-closure normalization is identified with the
subgroup-fixed normalization ring in
\normsrc{InvariantIntermediate.lean}.
\end{leanfiles}

\begin{remark}[Generic, not compactified]
Proposition~\ref{prop:diagonal-selector} is a generic algebra statement.  It
does not extend the selector across a compactification and does not identify
pointwise selector multiplication with graph convolution.
\end{remark}

\section{The marked-sheet block normal form}
\label{sec:marked-block}

Let \(G\) be a group, \(H\leq G\), and

\[
  \Sigma_H=G/H,
  \qquad F=F_{\Q}(H)=\Q^{G/H}.
\]

Finiteness of \(G/H\) is unnecessary for the linear identities in this
section.  Let \(s_0=H\) be the base sheet, let \(q_H\in F\) be its
characteristic function, and put

\[
  E_H=\ker(\ev_{s_0}:F\to\Q),
  \qquad
  a_H=q_H-\one\in E_H.
\]

There are two natural splittings of evaluation.  The selector splitting is

\[
  \Phi_{q_H}:F\xrightarrow{\sim}\Q\oplus E_H,
  \qquad
  \Phi_{q_H}(f)=
    \bigl(f(s_0),f-f(s_0)q_H\bigr),
\]

with inverse \((t,x)\mapsto tq_H+x\).  The constant-unit splitting is

\[
  \Phi_1:F\xrightarrow{\sim}\Q\oplus E_H,
  \qquad
  \Phi_1(f)=
    \bigl(f(s_0),f-f(s_0)\one\bigr),
\]

with inverse \((t,x)\mapsto t\one+x\).

Let \(m_{q_H}:F\to F\) denote pointwise multiplication by \(q_H\), and define

\[
  P_0(t,x)=(t,0),
  \qquad
  P_{a_H}(t,x)=(t,ta_H),
  \qquad
  U_{a_H}(t,x)=(t,x+ta_H).
\]

\begin{theorem}[Marked-sheet block normal form]
\label{thm:marked-block}
The marked selector satisfies

\begin{align}
  \Phi_{q_H}m_{q_H}\Phi_{q_H}^{-1}&=P_0,\label{eq:selector-diagonal}\\
  \Phi_1m_{q_H}\Phi_1^{-1}&=P_{a_H},\label{eq:selector-graph}\\
  \Phi_1\Phi_{q_H}^{-1}&=U_{a_H},\label{eq:selector-shear}\\
  P_{a_H}&=U_{a_H}P_0U_{a_H}^{-1}.\label{eq:shear-conjugacy}
\end{align}

In particular, \(m_{q_H}\), \(P_0\), and \(P_{a_H}\) are idempotent.
\end{theorem}

\begin{proof}
Take \((t,x)\in\Q\oplus E_H\).  Since \(x(s_0)=0\), pointwise multiplication
by \(q_H\) kills \(x\).  Therefore

\[
  m_{q_H}(tq_H+x)=tq_H,
\]

which proves \eqref{eq:selector-diagonal}.  In the constant-unit splitting,

\[
  m_{q_H}(t\one+x)=tq_H
  =t\one+t(q_H-\one),
\]

which proves \eqref{eq:selector-graph}.  Comparing the two inverse splittings
gives

\[
  \Phi_1\Phi_{q_H}^{-1}(t,x)
  =\Phi_1(tq_H+x)=(t,x+ta_H),
\]

and hence \eqref{eq:selector-shear}.  Conjugating
\eqref{eq:selector-diagonal} by this transition gives
\eqref{eq:shear-conjugacy}.  Idempotence also follows directly from
\(q_H^2=q_H\).
\end{proof}

\begin{leanfiles}
The two splittings, multiplication operator, graph operator, shear, and all
four identities are in \sheetsrc{MarkedBlockProjector.lean}.  The principal
theorems are \leanpath{decomposition_conj_multiplication},
\leanpath{selectorDecomposition_conj_multiplication},
\leanpath{selectorToMarked_eq_shear}, and
\leanpath{shear_conj_diagonalOperator}.
\end{leanfiles}

The same graph operator occurs in the Tate-block encoding of a rational
class.  If \(V\) is a rational vector space and \(\alpha\in V\), put

\[
  P_\alpha:\Q\oplus V\longrightarrow\Q\oplus V,
  \qquad
  P_\alpha(t,v)=(t,t\alpha).
\]

Evaluation of \(P_\alpha\) at \((1,0)\), followed by projection to \(V\),
recovers \(\alpha\).  Thus \(\alpha\mapsto P_\alpha\) is injective.

\begin{proposition}[Naturality of the marked graph]
\label{prop:marked-naturality}
For every rational linear map \(\theta:E_H\to V\), let

\[
  \Theta=\operatorname{id}_{\Q}\oplus\theta,
  \qquad
  \alpha=\theta(a_H).
\]

Then

\begin{equation}
\label{eq:graph-naturality}
  \Theta P_{a_H}=P_\alpha\Theta.
\end{equation}
\end{proposition}

\begin{proof}
Both sides send \((t,x)\) to \((t,t\theta(a_H))\).
\end{proof}

\begin{leanfiles}
The rational graph projector, its idempotence, shear form, and class
detection are in \chowsrc{TateGraphProjector.lean}.  The identification of
the marked-sheet graph operator with that construction and
Equation~\eqref{eq:graph-naturality} are in
\chowsrc{GaloisMarkedBlock.lean}, under
\leanpath{graphOperator_eq_classGraphProjector} and
\leanpath{productRealization_intertwines_selectorMultiplication}.
\end{leanfiles}

\begin{remark}[Naturality is not realization]
Proposition~\ref{prop:marked-naturality} applies to every rational linear map
\(\theta\).  It does not construct \(\theta\) from geometry, show that it is
motivic, or prove that \(\theta(a_H)\) is the class of an algebraic cycle.
\end{remark}

\section{A conditional realization criterion}
\label{sec:conditional-realization}

The exact remaining comparison can be stated without identifying the
pointwise sheet algebra with a correspondence algebra.  Let \(\Gbar\) be a
possibly noncommutative rational algebra.  Supply an
\leanpath{AlgebraProjectorExtension} \(\mathcal E\) with generic algebra
\(F\), compactified algebra \(\Gbar\), and generic projector \(q_H\).
Thus \(\mathcal E\) includes an algebra homomorphism

\[
  r:\Gbar\longrightarrow F,
\]

a type of extension witnesses, and for each witness \(x\) an idempotent
\(\mathcal E.\operatorname{projector}(x)\in\Gbar\) whose restriction is
exactly \(q_H\).  Choose such an \(x\), put
\(\Pi=\mathcal E.\operatorname{projector}(x)\), and let \(\Gbar\) act on
\(\Q\oplus V\):

\[
  \rho:\Gbar\longrightarrow
  \End_{\Q}(\Q\oplus V).
\]

Let \(\theta:E_H\to V\) be linear, put
\(\Theta=\operatorname{id}_{\Q}\oplus\theta\), and assume

\begin{equation}
\label{eq:action-square}
  \rho(\Pi)\Theta=\Theta P_{a_H}.
\end{equation}

This equation is substantive input: it compares the action of a compactified
correspondence with generic selector multiplication.

\begin{proposition}[Marked-selector realization criterion]
\label{prop:conditional-realization}
Assume Equation~\eqref{eq:action-square} and
\(\theta(a_H)=\alpha\in V\).  Then the lower-left class of \(\rho(\Pi)\) is
\(\alpha\):

\[
  \operatorname{pr}_V\bigl(\rho(\Pi)(1,0)\bigr)=\alpha.
\]

If \(\theta\) is surjective, then the whole operator is the graph projector,

\[
  \rho(\Pi)=P_\alpha.
\]
\end{proposition}

\begin{proof}
Evaluate \eqref{eq:action-square} at \((1,0)\).  Its right side is
\(\Theta(1,a_H)=(1,\alpha)\), which proves the first assertion.  If \(\theta\)
is surjective, then \(\Theta\) is surjective.  Proposition
~\ref{prop:marked-naturality} and \eqref{eq:action-square} give

\[
  \rho(\Pi)\Theta=\Theta P_{a_H}=P_\alpha\Theta.
\]

Surjective cancellation proves \(\rho(\Pi)=P_\alpha\).
\end{proof}

Now let \(M\) be a supplied rational cycle-class model, consisting of a
rational module \(\mathsf{Ch}\), a rational vector space \(V\), an independently
distinguished Hodge-class subspace of \(V\), and a cycle-class map

\[
  \operatorname{cl}:\mathsf{Ch}\longrightarrow V.
\]

The model also requires
\(\operatorname{im}(\operatorname{cl})\) to lie in that distinguished
subspace; it does not define the Hodge subspace to be the cycle-class image.

A lower-left detection map is a supplied linear map

\[
  d:\Gbar\longrightarrow\mathsf{Ch}
\]

compatible with the action in the sense that

\[
  \operatorname{cl}(d(z))
  =\operatorname{pr}_V\bigl(\rho(z)(1,0)\bigr)
  \qquad(z\in\Gbar).
\]

\begin{corollary}[Conditional cycle-class detection]
\label{cor:conditional-detection}
Under the hypotheses of Proposition~\ref{prop:conditional-realization} and a
supplied compatible lower-left detection map,

\[
  \operatorname{cl}(d(\Pi))=\alpha.
\]

Thus \(\alpha\) belongs to the range of the supplied cycle-class map.
Surjectivity of \(\theta\) is not needed for this conclusion.
\end{corollary}

\begin{leanfiles}
The comparison record, lower-left calculation, surjective cancellation, and
cycle-class conclusion are in \chowsrc{GaloisMarkedRealization.lean}, under
\leanpath{MarkedGaloisSelectorComparison},
\leanpath{lowerLeftClass_action_eq}, \leanpath{realizes}, and
\leanpath{isAlgebraicHodgeClass_of_comparison}.  The cycle-class model and
the definition of its cycle-class range are in
\leanpath{HodgeConjecture/Realization/HodgeClasses.lean}.
\end{leanfiles}

\begin{remark}[Why the ambient algebra may be noncommutative]
The graph projectors satisfy

\[
  P_\alpha P_\beta=P_\alpha,
  \qquad
  P_\beta P_\alpha=P_\beta.
\]

They commute exactly when \(\alpha=\beta\).  Therefore one fixed commutative
correspondence algebra acting through an algebra homomorphism cannot realize
distinct graph projectors.  The conditional interface deliberately keeps
the generic sheet algebra commutative while permitting \(\Gbar\) to be
noncommutative.
\end{remark}

\begin{remark}[No Hodge consequence in the present draft]
Corollary~\ref{cor:conditional-detection} is a theorem about supplied abstract
data.  The repository does not construct, for a smooth projective complex
variety, the Chow module \(\mathsf{Ch}\), Betti cohomology \(V\), the Hodge-class
subspace, the Tate-augmented correspondence algebra, the compactified
extension package \(\mathcal E\), its restriction map or idempotent \(\Pi\),
the residual map \(\theta\), Equation
~\eqref{eq:action-square}, or the detection map \(d\).  In particular, it
does not establish a new class of varieties or Hodge classes satisfying the
Hodge conjecture.
\end{remark}

\subsection{A transparent conditional boundary corollary}
\label{subsec:conditional-boundary}

The formal development also records the following forward-only endpoint.
Let \(R\to A\) be an essentially finite-type map of commutative rings, let
\((\mathfrak p_\lambda)_{\lambda\in\Lambda}\) be a chosen family of prime ideals of
\(A\), and let \((I_\lambda)_{\lambda\in\Lambda}\) be a chosen family of
inertia subgroups of a deck-transformation group \(G\).  Fix a finite subgroup
\(H\leq G\), and write \(\pi_H^\eta\) for the generic subgroup-average
correspondence.  Supply generic and compactified rational algebras
\(\Geta\) and
\(\Gbar\), a rational Chow-correspondence model \(\chi\),
and a prescribed rational cohomological projector with underlying operator
\(T\).  Write \(\operatorname{cl}_{\chi,\mathrm{gen}}\),
\(\operatorname{res}_{\chi}\), and \(\rho_{\chi}\) for the model's generic
cycle-class, restriction, and cohomological-action maps.

\begin{corollary}[Conditional different-avoidance boundary endpoint]
\label{cor:conditional-boundary}
Suppose the following data and implications are supplied separately:

\begin{enumerate}[label=\textup{(\roman*)},leftmargin=2.6em]
  \item a type \(\mathcal X\) of proposed extensions and a compactified-class
        extraction map \(c:\mathcal X\to\Gbar\);
  \item compactified idempotence, \(c(x)^2=c(x)\) for every
        \(x\in\mathcal X\);
  \item the restriction identity
        \[
          \operatorname{res}_{\chi}(c(x))
          =\operatorname{cl}_{\chi,\mathrm{gen}}(\pi_H^\eta)
          \qquad(x\in\mathcal X);
        \]
  \item the one-way purity implication that simultaneous inertia saturation
        of the base sheet under every \(I_\lambda\) yields a witness in
        \(\mathcal X\);
  \item for every chosen prime \(\mathfrak p_\lambda\), a supplied
        \leanpath{HeightOneInertiaRealization} comparing the relevant
        algebraic ramification data with \(I_\lambda\), \(H\), and
        \(\mathfrak p_\lambda\);
  \item the prescribed-action identity
        \[
          \rho_{\chi}(c(x))=T
          \qquad(x\in\mathcal X);
        \]
  \item avoidance of the collision Noether different at every chosen prime:
        \[
          \mathfrak D_{A/R}^{\mathrm{coll}}\not\subseteq \mathfrak p_\lambda
          \qquad(\lambda\in\Lambda).
        \]
\end{enumerate}

Then there exists \(\Pi\in\Gbar\) such that

\[
  \Pi^2=\Pi,
  \qquad
  \operatorname{res}_{\chi}(\Pi)
    =\operatorname{cl}_{\chi,\mathrm{gen}}(\pi_H^\eta),
  \qquad
  \rho_{\chi}(\Pi)=T.
\]
\end{corollary}

\begin{proof}
Essential finite type supplies collision-different detection at every
\(\mathfrak p_\lambda\).  For each \(\lambda\), combine that detection, the supplied
\leanpath{HeightOneInertiaRealization}, and different avoidance with the
base-sheet extension criterion to obtain \(I_\lambda\)-saturation of the
base sheet.  The one-way purity implication then gives some
\(x\in\mathcal X\).
Taking \(\Pi=c(x)\), the separately supplied idempotence, restriction, and
prescribed-action hypotheses give the three displayed conclusions.
\end{proof}

The essential finite-type hypothesis supplies only the detection of the
non-unramified locus by the collision Noether different, in the sense of
Stacks Project, Lemma~49.6.8, Tag~\texttt{0BVU}
\cite{StacksNoetherDifferent}.  It does not supply the inertia realization,
purity implication, compactified class, restriction comparison, or
prescribed-action equality.  The primes \(\mathfrak p_\lambda\) are merely the chosen prime
ideals in the Lean statement; they are not encoded there as an exhaustive
boundary family or as height-one primes.  The name of the supplied inertia
interface records its intended height-one application.

\begin{leanfiles}
This is
\leanpath{exists_prescribed_subgroupAverage_chowProjector_of_different_away}
in \chowsrc{GaloisProjector.lean}.  Its conclusion is an idempotent in an
abstract compactified algebra with restriction and action equations, not a
constructed element of a concrete Chow ring.
\end{leanfiles}

\begin{remark}[Exact force of the boundary slogan]
Corollary~\ref{cor:conditional-boundary} is the honest current formal form of
``different avoidance extends the projector.''  It is a one-way conditional
packaging theorem, not the desired unconditional boundary-extension theorem
and not an equivalence.
\end{remark}

\section{Formal status and limitations}
\label{sec:status}

The verified dependency spine is

\[
\begin{gathered}
  \text{finite-\'etale span classes}
  \longrightarrow
  \mathcal C_{\fe}^{\mathrm{free}}(X/S),\\
  \text{deck action}
  \longrightarrow
  \Q[G]\longrightarrow
  \mathcal C_{\fe}^{\mathrm{free}}(X/S),\\
  \text{group/Hecke idempotent}
  \longrightarrow
  \text{formal convolution projector},\\
  \text{supplied compactified graph square}
  \longrightarrow
  \text{compactified idempotent with correct restriction},\\
  \text{marked pointwise selector}
  \longrightarrow
  \text{diagonal and graph block forms}.
\end{gathered}
\]

The conditional Galois collision theorem additionally consumes the explicit
torsor/coproduct comparison and a finite-additive algebra realization.  The
conditional marked realization criterion additionally consumes the
compactified extension, residual map, action square, and detection map.  The
conditional boundary corollary consumes the separately listed class
extraction, compactified idempotence, purity, inertia, restriction, and
prescribed-action inputs; essential finite type discharges only different
detection.  The following are not results of this note:

\begin{itemize}[leftmargin=2.2em]
  \item identification of the free span algebra with a Chow or bivariant
        correspondence algebra;
  \item extraction of pure-dimensional proper cycle classes or compatibility
        with refined Chow convolution;
  \item extension of a generic selector or cover across a compactification;
  \item equivalence of selector invariance, inertia triviality, and vanishing
        of a different or conductor obstruction;
  \item construction of a Chow-motivic residual map or proof of the action
        square \eqref{eq:action-square};
  \item a concrete Chow, cohomology, Hodge-structure, or cycle-class model;
  \item any partial or complete proof of the Hodge conjecture.
\end{itemize}

In particular, a pointwise idempotent in a sheet-function algebra must not be
advertised as a Chow idempotent.  Even when a generic cycle has a closure,
compatibility of restriction with correspondence composition shows only that
its idempotence defect lies in the restriction kernel.  Identifying that
kernel geometrically with boundary-supported classes, and then proving that
the particular defect vanishes or can be corrected, require additional
arguments.

The checked development uses Lean~4.24.0 and Mathlib revision
\texttt{f897ebcf72cd16f89ab4}\allowbreak
\texttt{577d0c826cd14afaafc7}.  Lean and Mathlib are described in
\cite{Lean4,Mathlib}.  An immutable release commit or archive has not yet
been assigned.

The working source is at
\url{https://github.com/what-social-construct/HodgeConjecture}.  Before any
submission, the statements and source-line mappings must be checked against
the release commit, the bibliography and neighboring literature must be
audited, author and license metadata must be finalized, and an independent
mathematical review must be recorded.

\section{Conclusion}
\label{sec:conclusion}

Finite-\'etale scheme spans provide a small admissible correspondence class
whose convolution is stable under the basic geometric operations used by
finite group actions.  Passing to isomorphism classes and then to the rational
monoid algebra makes the convolution laws strict.  Graphs consequently turn
every rational group-algebra idempotent into a formal correspondence
idempotent, transport orthogonal group-algebra decompositions to decompositions
of the diagonal, produce normalized finite-group projectors and
Hecke-corner/double-coset operators, and transport supplied idempotents of the
Hecke corner.

The marked-sheet calculation supplies a separate and exact bridge at the
linear level.  A base-sheet selector is diagonal in selector coordinates but
becomes the canonical graph projector of \(a_H=q_H-1\) in constant-unit coordinates.
This is a change of splitting, not an identification of pointwise and
convolution products.  Naturality under a residual linear map remains pure
linear algebra.

The conditional realization criterion shows precisely what further data
would connect that block identity to a cycle-class statement.  Constructing
those data geometrically remains open in this repository.  The present
contribution is therefore the formal correspondence calculus, the supplied
compactified-graph comparison and raw-defect calculus, the marked block
theorem, and auditable interfaces for the missing realization and
boundary-extension steps.

\paragraph{Acknowledgment.}
OpenAI Codex assisted with the Lean formalization and manuscript draft.  The
manuscript remains subject to author review; responsibility for its final
statements, proofs, citations, and release terms belongs to the author.

\begingroup
\small
\bibliographystyle{alpha}
\bibliography{references}
\endgroup

\end{document}
